\ulohahead{společná informatika}{Synchronizace: kritická sekce, semafory, producent-konzument (C\#)}
\begin{zadani}
\begin{enumerate}[label=(\alph*)]
  \item \bod{1} Definujte časově závislou chybu (race condition), kritickou sekci a vzájemné vyloučení. Co je zámek a jaký je rozdíl mezi aktivním a pasivním čekáním?
  \item \bod{2} Co je semafor (operace \textsc{wait}/P a \textsc{signal}/V)? Napište v C\# řešení producent-konzument s omezeným bufferem pomocí semaforu.
  \item \bod{3} Uveďte čtyři podmínky uváznutí (deadlock). Porovnejte spinlock a mutex a vysvětlete atomickou operaci compare-and-swap (CAS).
\end{enumerate}
\end{zadani}

\begin{reseni}
\textbf{(a)} \emph{Race condition}: výsledek závisí na nedeterministickém prokládání vláken nad sdíleným stavem. \emph{Kritická sekce}: úsek kódu, který smí běžet nejvýše jedno vlákno najednou. \emph{Vzájemné vyloučení} (mutual exclusion) tuto vlastnost zajišťuje. \emph{Zámek} (lock) chrání kritickou sekci; \emph{aktivní čekání} (spin) cyklí a pálí CPU, \emph{pasivní} vlákno uspí a plánovač ho probudí (lepší při delším čekání).

\textbf{(b)} \emph{Semafor} = čítač s atomickými operacemi: \textsc{wait}/P sníží čítač (a zablokuje při 0), \textsc{signal}/V zvýší (a probudí čekající). Omezený buffer kapacity $N$:
\begin{lstlisting}
// empty: kolik je volnych mist, full: kolik je obsazenych
var empty = new SemaphoreSlim(N, N);
var full  = new SemaphoreSlim(0, N);
var mutex = new object();
var buffer = new Queue<int>();

void Producer(int item) {
    empty.Wait();                 // pockej na volne misto (P)
    lock (mutex) { buffer.Enqueue(item); }
    full.Release();               // pribyl prvek (V)
}

int Consumer() {
    full.Wait();                  // pockej na prvek (P)
    int item;
    lock (mutex) { item = buffer.Dequeue(); }
    empty.Release();              // uvolnil se slot (V)
    return item;
}
\end{lstlisting}

\textbf{(c)} \emph{Deadlock} nastane při současném splnění: (1) vzájemné vyloučení, (2) drž a čekej, (3) bez odebrání (no preemption), (4) kruhové čekání. \emph{Spinlock} aktivně čeká (vhodný pro velmi krátké sekce na více jádrech), \emph{mutex} uspí vlákno (vhodný pro delší). \emph{CAS}\,(\texttt{CompareAndSwap}(addr, expected, new)): atomicky porovná a případně zapíše; základ neblokujících (lock-free) struktur.
\end{reseni}

\kalibrace{Synchronizace (semafory, kritická sekce, race condition) je nejčastější OS podtéma a často chce \emph{kód}. Definice (a) = jistý bod, kostra semaforu (b) = druhý; tím se vyhneš nule na implementační otázce.}
\past{Pořadí semaforu u producent-konzument je klíčové: čekej na zdroj (\textsc{wait}) PŘED zámkem, jinak hrozí deadlock. \texttt{empty} a \texttt{full} se nesmí prohodit.}
